%------------------------------------------------------
% Harvard-style one-page resume -- NGUYEN KHAC BAO
% Target: AI Engineer Intern / RAG & GenAI
% Compile: pdflatex CV_Nguyen_Khac_Bao_AI_Engineer_Intern_TOEIC.tex
%------------------------------------------------------

\documentclass[10pt,a4paper]{article}

\usepackage[T1]{fontenc}
\usepackage[utf8]{inputenc}
\usepackage[a4paper,top=0.42in,bottom=0.42in,left=0.55in,right=0.55in]{geometry}
\usepackage{enumitem}
\usepackage{hyperref}
\usepackage{titlesec}
\usepackage{microtype}
\usepackage{xcolor}

\hypersetup{
    colorlinks=true,
    urlcolor=blue,
    linkcolor=black,
    pdfborder={0 0 0}
}

\titleformat{\section}
    {\normalfont\large\bfseries}
    {}
    {0em}
    {}
    [\vspace{1pt}\titlerule]

\titlespacing*{\section}{0pt}{5pt}{2pt}

\setlength{\parindent}{0pt}
\setlength{\parskip}{0pt}
\pagestyle{empty}
\sloppy

\newenvironment{cvitems}{%
    \begin{itemize}[leftmargin=1.25em,itemsep=0pt,topsep=1pt,parsep=0pt]
}{%
    \end{itemize}
}

\newcommand{\entryhead}[2]{%
    \noindent\textbf{#1}\hfill #2\\[-2pt]
}

\newcommand{\entrysub}[1]{%
    \noindent\textit{#1}\\[-1pt]
}

\begin{document}
\fontsize{9.35}{10.35}\selectfont

\begin{center}
    {\Large\textbf{NGUYEN KHAC BAO}}\\[-1pt]
    AI Engineer Intern\\[2pt]
    Ho Chi Minh City, Vietnam
    \,\textbar\,
    \href{mailto:nguyenkhacbaowork564@gmail.com}{nguyenkhacbaowork564@gmail.com}
    \,\textbar\,
    (+84) 039 586 4429
    \,\textbar\,
    \href{https://github.com/NguyenKhacBao564}{github.com/NguyenKhacBao564}
\end{center}

\section{CAREER OBJECTIVE}
AI Engineer Intern candidate with hands-on experience building RAG/GenAI systems, fine-tuning LLMs, and developing backend data pipelines. Seeking to apply a Software Engineering mindset to build trustworthy AI applications for complex domains.

\section{EDUCATION}
\entryhead{Posts and Telecommunications Institute of Technology (PTIT)}{2022 -- 2027}
\entrysub{Bachelor of Information Technology}
Relevant Coursework: Machine Learning, Deep Learning, Computer Vision.

\section{TECHNICAL SKILLS}
\textbf{Programming Languages:} Python, SQL (basic), Bash\\
\textbf{AI / Machine Learning:} RAG pipelines, LlamaIndex ReAct, PEFT/LoRA, Prompt Engineering, PyTorch\\
\textbf{Computer Vision:} YOLOv8/v11, YOLO-Pose, OpenCV, ByteTrack, HyperLPR3/OCR\\
\textbf{Tools \& Backend:} FastAPI, Docker, Qdrant (Vector DB), Git/GitHub, Linux\\
\textbf{Basic Cloud \& Messaging:} Vast.ai (GPU), AWS EC2, Google Cloud Run (microservices), RabbitMQ

\section{LANGUAGES}
\entryhead{TOEIC: 610/990}{2026}
\entryhead{English: Technical reading proficiency}{}
Able to read official documentation, research papers, model cards, API references, and engineering articles in English.

\section{PROJECTS}
\entryhead{Vietnamese Legal RAG Chatbot}{2026}
\entrysub{Personal Project \textbar\ Technology: Python, FastAPI, Streamlit, Celery, Valkey, MariaDB, Qdrant, PEFT/LoRA, Llama 3.1, llama.cpp \textbar\ Source: \href{https://github.com/NguyenKhacBao564/legal_RaG_Chatbot_VietNamese-System}{GitHub}}
\begin{cvitems}
    \item Built an end-to-end Vietnamese legal assistant with a FastAPI backend, Streamlit chat UI, Celery asynchronous task handling, Valkey/Redis queueing, and MariaDB conversation storage.
    \item Designed an agent-style RAG workflow to route user queries between legal retrieval, LlamaIndex ReAct tools, web-search handling, and chat fallback.
    \item Implemented hybrid retrieval over Qdrant vector database, multi-query expansion, reranking, and legal-context prompting for Vietnamese legal QA.
    \item Fine-tuned Llama-3.1-8B-Instruct with PEFT/LoRA on Vietnamese legal QA data using Vast.ai; served the adapter through an OpenAI-compatible local API.
    \item Optimized local inference by converting the serving path to llama.cpp with GGUF Q4, reducing memory pressure compared with full FP16 Transformers inference.
\end{cvitems}

\entryhead{FallGuard Labeling Tool -- Computer Vision Dataset Assistant}{2026}
\entrysub{Personal Project \textbar\ Technology: Python, Streamlit, OpenCV, Pandas, YOLO-Pose, PyYAML, pytest \textbar\ Source: \href{https://github.com/NguyenKhacBao564/autofall-labeler-tool}{GitHub}}
\begin{cvitems}
    \item Built a Python-based dataset assistant for fall-detection image/video datasets, covering raw data import, video frame extraction, metadata generation, human review, QA checks, split, and export.
    \item Integrated YOLO-Pose to generate 13-point human pose annotations and built a Streamlit UI for label correction, uncertain/excluded samples, pose review, and manual keypoint editing.
    \item Designed reproducible CSV source-of-truth records using metadata, annotations, keypoints, and annotation events; tested with Kaggle Fall Detection and Multiple Cameras Fall datasets.
\end{cvitems}

\entryhead{Traffic Violation Detection \& License Plate Recognition}{2026}
\entrysub{Personal Project \textbar\ Technology: Python, OpenCV, Ultralytics YOLO, ByteTrack, HyperLPR3, NumPy, Shapely, JSON \textbar\ Source: \href{https://github.com/NguyenKhacBao564/Traffic-Violation-Detection-and-License-Plate-Recognition}{GitHub}}
\begin{cvitems}
    \item Built a fixed-camera traffic video pipeline for vehicle detection, multi-object tracking, traffic-light state estimation, stop-line violation logic, license plate detection, OCR, and evidence generation.
    \item Fine-tuned a YOLO-based license plate detector on a compact CCPD2019 subset of 8,598 images, achieving 0.989 precision, 0.998 recall, and 0.994 mAP50 on validation.
    \item Implemented dataset QA for broken files, YOLO bbox validity, negative sample logic, OCR crop availability, duplicate review, low-light images, heavy blur, and small license plates.
\end{cvitems}

\section{ACTIVITIES}
\entryhead{AI Research Group -- University Club}{2025 -- Present}
\entrysub{Member}
Discussed Machine Learning, Deep Learning, Computer Vision, and AI application development.

\end{document}
